\documentclass[a4paper,11pt]{article}

\usepackage{revy}
\usepackage[utf8]{inputenc}
\usepackage[T1]{fontenc}
\usepackage[danish]{babel}
\usepackage{amsmath,amssymb,amsfonts,amsthm,textcomp}


\revyname{MatematikRevyen}
\revyyear{2026}
% HUSK AT OPDATERE VERSIONSNUMMER
\version{0.1}
\eta{$2$ minutter}
\status{Færdig}

\title{Vinsmagning}
\author{Carl '21}

\begin{document}
\maketitle

\begin{roles}
\role{S1}[Skuespiller] Studerende
\role{S2}[Skuespiller] Studerende
\role{S3}[Skuespiller] Sommelier
\end{roles}

\begin{props}
\item Bord med to stole, fancy stil
\item 2 Menukort
\item 3 flasker vin
\end{props}


\begin{sketch}
\scene{S1 og S2 er fint klædt på og sidder ved et bord på scenen. S3 kommer ind og giver nogle menukort til S1 og S2}

\says{S3} Velkommen, hvad kunne I tænke jer at smage i dag?

\says{S1} Jeg kunne godt tænke mig at starte med noget simpelt.

\says{S2} Bare en klassisk konvergenssætning til mig, tak!

\says{S3} Så gerne

\scene{S3 tager nogle flaske frem}

\says{S3}[til S1] Her er en gammel kending, "$\sqrt{2}$". Den kan drikkes af de fleste. Bevist ved modstrid, så smager den bedste. Prøv at smag

\scene{S3 hælder en sjat op i S1's glas. S1 smager og spytter ud}

\says{S1} Føj! Jeg tror der er gået induktion i den, det gør altså ikke noget godt for beviset.

\says{S3}[undskyldende] Det er jeg frygtelig ked af. Den må være blevet vejledt forkert. Jeg henter straks en ny!

\says{S1} Nej, nej. Jeg har skiftet mening. Kom  med noget moderne i stedet

\says{S3} Det får du.
\says{S3}[til S2] Til dig har jeg hentet en personlig favorit, "Store Tals Lov". Den fås i utallige varianter med forskellige antagelser. Se om du kan smage hvilken en det her er

\scene{S3 hælder op. S2 dufter til den og smager på den}

\says{S2} Mmmh, den er fyldt med gode idéer. Borel-Cantelli, Kolmogorov. Lidt ergodicitet, men ikke for meget. Uhh, er det et endeligt andet moments jeg smager? Nej! Endeligt første moment!!

\says{S3}[nikker] Imponerende, du har en skarp intuition.

\says{S3}[til S1] Jeg er straks tilbage med en ny flaske til dig

\scene{S3 forlader scenen}

\says{S1}[til S2] Må jeg lige smage din?

\says{S2}[giver glasset] Selvfølgelig

\says{S1} Uh nej, det er ikke min smag det der sandsynlighedsteori. Det er som som man aldrig er \textit{helt} sikker på hvad man får

\says{S2} Sandt, men jeg har næsten altid gode oplevelser med det alligevel

\scene{S3 kommer tilbage med en flaske}

\says{S3}[til S1] Det her er vores mest moderne. Den har været under produktion længe og er lige blevet klar. Jeg har ikke selv fået smagt den endnu desværre. Den hedder "Navier-Stokes"

\scene{S3 hælder op og S1 smager på den}

\says{S1} Mmm, meget imponerende. Dynamisk, rummelig. Glat og lækker. Virkelig en vin i verdensklasse.

\says{S1}[skifter udtryk] Eeev!! Hvad er det for en finish?!? Det efterlader godt nok en med en dårlig smag i munden...

\says{S3}[læser på flasken] "Kan forekomme singulariteter". Hmm, måske den ikke var helt klar til at blive drukket endnu.

\scene{Tæppe}


\end{sketch}
\end{document}
